\appendix

\section{Fluorescence and Electron-Microscopy Fidelity Paths}
\label{sec:appendix-non-scalar-modalities}

\paragraph{Transmission electron microscopy phase contrast.} The transmission-electron-microscopy path supports configurable backend fidelity through \texttt{tem\_backend}: weak-phase CTF proxy (\texttt{ctf\_proxy}), reduced-slice proxy (\texttt{multislice\_lite}), or Syniscopy split-step multislice (\texttt{syniscopy\_multislice}). It exposes accelerating voltage, Cs, defocus, partial coherence, slice count/thickness, and objective-aperture controls, with CTF and reduced-slice modes kept as lightweight diagnostics and syniscopy\_multislice used for transport-aware comparisons when explicitly selected.

\paragraph{Scanning electron microscopy secondary-electron contrast.} The scanning-electron-microscopy path computes a probe-blurred mixture of particle bulk signal, particle edge signal, mounting-interface topography, and material-dependent secondary-electron yield. The default branch uses \texttt{monte\_carlo\_transport} with a sliced SEM source volume, stochastic interaction-volume kernels, explicit beam-energy, dwell-time, beam-current, and detector-geometry bookkeeping; alternate branches are \texttt{syniscopy\_transport\_lite}, \texttt{gaussian\_probe\_proxy}, \texttt{interaction\_volume\_proxy}, and \texttt{reference\_kernel\_table} interpolation paths. Interaction-volume coupling, topography coupling, detector acceptance, escape-depth, and charge-response behavior are configurable terms in the transport-calibrated paths, with simpler proxy terms retained when a compact option is selected.

\paragraph{Widefield and total internal reflection fluorescence.} The fluorescence paths integrate \path{MaterialProperties.fluorophore_density} over the particle footprint, modulate by excitation/emission spectral factors derived from \path{fluorescence\_excitation\_wavelength\_nm} and \path{fluorescence\_emission\_wavelength\_nm}, convolve with a vectorial-emission point-spread-function backend, and add sample-environment autofluorescence/excitation modulation. The fluorescence response is determined by the material and particle configuration supplied to the simulator. The \texttt{vectorial\_photophysics} branch adds frame-wise Markov-state occupancy (blinking/recovery), bleaching, detector-quantum-efficiency scaling, and optional TIRF depth attenuation; the parametric PSF branch remains the lightweight baseline.

\section{Bench-to-Simulator Parameter Mapping}
\label{sec:bench-mapping}

This appendix is a configuration reference. The physics is stated in Section~\ref{sec:physical-model}; Table~\ref{tab:bench-mapping} records which Syniscopy configuration parameter carries each measured bench quantity, so that measured bench quantities can be represented in, or used to seed, a simulator configuration.

\begin{table}[H]
\centering
\small
\begin{tabular}{p{0.30\linewidth}p{0.14\linewidth}p{0.46\linewidth}}
\toprule
Bench measurement & Units & Simulator parameter \\
\midrule
Illumination wavelength & nm & \texttt{wavelength\_nm} \\
Electron accelerating voltage & keV & \texttt{tem\_acceleration\_kV} \\
Objective numerical aperture & dimensionless & \texttt{numerical\_aperture} \\
Effective pixel pitch at sample plane & nm & \texttt{pixel\_size\_nm} \\
Detector read noise & camera counts RMS & \texttt{read\_noise\_counts} \\
Detector quantum efficiency & dimensionless & \texttt{detector\_qe} and modality-specific aliases; when \texttt{detector\_input\_is\_incident\_quanta} is true, QE converts incident quanta to detected quanta before shot-noise variance \\
Per-frame detected-count or dose scale & counts/electrons & \texttt{background\_intensity}, \texttt{fluorescence\_photon\_count\_scale}, \texttt{tem\_dose\_per\_pixel}, \texttt{sem\_electrons\_per\_pixel}, or analysis-level NQ for detected-quanta-normalized diagnostics \\
Frame rate & Hz & \texttt{fps} \\
Particle material or complex index & $n,k$ & \texttt{material}; \texttt{refractive\_index} \\
Nominal particle diameter & nm & \texttt{diameter\_nm} \\
Medium refractive index / viscosity & dimensionless; Pa\,s & \texttt{medium\_material}, \texttt{viscosity\_Pa\_s} \\
Mounting-interface pattern & enum & \texttt{sample\_environment\_pattern} \\
Hole or grid spacing & nm & \texttt{substrate\_pattern\_dimensions} or pattern-specific recipe keys \\
\bottomrule
\end{tabular}
\caption{Common bench measurements and the corresponding Syniscopy configuration keys.}
\label{tab:bench-mapping}
\end{table}

A cropped experimental recording with pixel pitch $p$ and crop size $L\times L$ maps to a simulator pixel size of $p$ and an image size of $L$ pixels; if the experimental recording was binned $2\times2$, use twice the native camera pitch. Syniscopy writes background-subtracted lossless PNG frame sequences for training and can optionally save raw signal/reference frame arrays for audit. Pattern fitting from experimental recordings is available for the implemented gold-hole pattern family, allowing pitch and radius estimates to seed a bench-matched sample-environment configuration while unresolved jitter and shape-regularity fields remain explicit metadata. A one-to-one pixel match to a real bench remains approximate because stray light, dust, gain nonlinearity, out-of-focus particles, and alignment errors are not fully absorbed by the current parameter set; source-matching residual structure is exposed through explicit sample-environment roughness/speckle fields \texttt{sample\_environment\_pattern\_roughness\_*} (static, flicker, or source-matched modes) with source coupling options \texttt{independent}, \texttt{coherent\_amplitude}, \texttt{field\_weighted}, \texttt{scene\_weighted}, and \texttt{channel\_weighted} (default: \texttt{channel\_weighted}).

The detected-quanta-normalized analyses use an analysis-level common budget $N_Q$ recorded in experiment output together with Contract-Q normalization metadata and derivative-convergence inheritance fields. That analysis budget is distinct from the default public configuration dictionary, where cross-modality comparisons use modality-specific count/dose scale parameters unless explicitly configured for a post-hoc normalization sweep.

\section{Literature Source-Use Audit}
\label{sec:source-use-audit}

The source-use audit in \Cref{tab:native-regime-reference-check} records which external source motivates each native-regime profile, what quantity that source reports, and whether the row is parameter-matched to the cited experiment. The source-use category distinguishes direct quoted localization precision, formula- or theory-derived localization scale, literature localization scale, modality-principle citation only, dimensional metrology scale, and rows without a quoted localization bound. Rows outside the localization-scale categories report \texttt{n/a} in the computed localization column, so dimensional-metrology and modality-principle citations are not visually treated as center-localization validations. The audit sources are carried as formal references with DOI metadata where available \cite{kovari-brightfield-3d,huang-cobri-virus,ando-darkfield-angstrom,juffmann-iscat-cobri-crlb,kurata-zernike-phase-retrieval,tian-waller-dpc,bon-qpi-nanoparticles,verpillat-dhm-gold,clack-groves-ricm,thompson-larson-webb,axelrod-tirf,rice-tem-sizing,crouzier-sem-diameter}.

\IfFileExists{figures/native-regime-localization-comparison-table.tex}{%
\begin{table}[H]
\centering
\begingroup
\setlength{\tabcolsep}{3pt}
\renewcommand{\arraystretch}{0.78}
\scriptsize
\resizebox{\linewidth}{!}{\input{figures/native-regime-localization-comparison-table}}
\endgroup
\caption{Native-regime source-use audit. The computed-localization column is populated only for source categories that report or derive a localization scale. Rows marked dimensional metrology, modality principle, or no quoted bound are carried for citation provenance and report \texttt{n/a} in that column. The parameter-match column records whether the source supplies enough optical, detector/dose, sample, and target-bound information for a quote-matched comparison.}
\label{tab:native-regime-reference-check}
\end{table}
}{}

\section{Reproducibility Workflow}
\label{sec:reproducibility}

The supplemental notebooks generate the raw experiment outputs used by the manuscript when executed in order. The paper-facing tables and figures are assembled from those outputs by \path{paper/assemble_output_artifacts.py}. That script writes \path{paper/figures/artifact-provenance-manifest.json}; rerunning E03 writes \path{supplemental/outputs/E03/comparison_contracts_manifest.json} and \path{supplemental/outputs/E03/artifact_dependency_graph.json}, so paper artifacts can be traced back to comparison contracts, source-output files, generator paths, and SHA-256 hashes. The results reported here were produced with the current checked-out Syniscopy notebook bundle in this repository. The deterministic simulator analyses use fixed seeds and run on a CPU runtime. Segment Anything Model 2 fine-tuning and inference write checkpoint, cache, prompt, and loss-semantics manifests so the GPU workflow can be reconstructed from recorded runtime and checkpoint metadata. Larger arrays, videos, checkpoints, and inference outputs can be regenerated from the notebooks.

Core comparison claims are contract-specific and map to explicit output files in \path{supplemental/outputs/} when the corresponding notebook pipeline is run:
\begin{itemize}
\item Contract-LP (lateral Cram\'{e}r--Rao ranking): when generated, \path{supplemental/outputs/E03/modality_ranking.csv} and \path{supplemental/outputs/E03/derivative_convergence_crlb.csv}.
\item Contract-LZ (\(\mathrm{SE}(3)\) and axial diagnostics): when generated, \path{supplemental/outputs/E03/cross_modality_axial.csv} and \path{supplemental/outputs/E03/axial_convergence_crlb.csv}; per-composite orientation support is in \path{supplemental/outputs/E02/orientation_crlb.csv} and assembled into \path{figures/orientation-crlb-table.tex}.
\item Contract-Q (detected-quanta normalized comparison): when generated, \path{supplemental/outputs/E03/detected_quanta_normalized.csv} and \path{supplemental/outputs/E03/detected_quanta_normalization_metadata.csv}, plus budget sweep artifacts in \path{figures/detected-quanta-budget-pareto.png}.
\item Contract-NR (native-regime source-use context): when generated, \path{figures/native-regime-localization-comparison-table.tex}, produced from the source-use audit rows, and associated rows in \path{supplemental/data/liverpool_caustic_50nm_review/}.
\end{itemize}
The deterministic simulator-side checks and notebook diagnostics in Sections~\ref{sec:results} and \ref{sec:appendix-non-scalar-modalities} are intended for conceptual and software validation; they are not substitutes for physical-instrument calibration.

The dataset-generation routine seeds each corpus from the caller-provided seed plus the video index. The dataset metadata records per-video parameter dictionaries, the supervision policy, and the annotation schema, enabling regeneration of individual videos.

The caustic-video sources used by the transfer-inference notebook are retained review-manifest rows derived from the University of Liverpool DataCat dataset by Schleyer \cite{schleyer-caustic}, licensed Creative Commons Attribution 4.0. The local \path{50nm/} DataCat download is a working input and is not distributed as Syniscopy source code. Reviewed clip windows, prompt records, and per-clip audit metadata are exported under \path{supplemental/data/liverpool_caustic_50nm_review/}; the manifest records the source identifier, selected source-frame window, prompt frame, prompt coordinates, 40~fps acquisition timebase, and 0.116~$\mu$m pixel size used by the transfer analysis. Reviewed clips derived from DataCat videos retain the DataCat attribution notice. Generated transfer overlays and masks are reproducible from the notebooks and are not required as public archival inputs. The Segment Anything Model 2 notebooks use Meta Segment Anything Model 2 code and checkpoints; those external assets and any fine-tuned checkpoint derived from them retain Meta's upstream licenses and notices and are not relicensed as Syniscopy code.

\section{Segment Anything Model 2 Dataset Interface}
\label{sec:sam2-dataset-interface}

The Segment Anything Model 2 example uses Syniscopy datasets as input to a downstream video-segmentation model. The training examples use lossless background-subtracted PNG frame sequences and target masks; preview videos are illustrative only. Optional raw signal/reference arrays serve as diagnostic records. The Segment Anything Model 2 workflow selects the requested target mask and propagates ignore masks and per-pixel loss weights so ignored pixels contribute zero loss.


\paragraph{Source-use separation.}
Native-regime source use is split into two generated tables. The localization-comparison table contains computed Syniscopy CRLB rows only when the cited source reports a localization scale. The source-provenance table contains citation and parameter-use cards only; those rows do not carry computed CRLB values and are not validation rows.
\input{figures/native-regime-source-provenance-table}


\paragraph{Convergence and artifact status.}
The paper-facing E03 tables are gated by generated convergence/status fields. Contract-LP rankings use selected rerendered finite-difference lateral Fisher matrices from \texttt{derivative\_convergence\_crlb.csv}. Contract-Q rankings use a separate detected-quanta derivative-convergence artifact rather than inheriting the lateral gate. Rows whose status is not validated remain diagnostic records.


\paragraph{SAM2 scope.}
The SAM2 notebooks and manifests describe a caustic-matched worked example. The E06 manifest records retained-clip accounting without inferring an unavailable candidate-pool size, E07/E08 record the caustic-matched training scope, and E09 persists per-clip mask-centroid motion metrics for the descriptive motion table.
